\section{Structure of solutions with respect to a packing}\label{sec:packing}
This section collects counting lemmas that hold for every packing $H$ (no maximality unless stated) and describe what a solution of size at most $|H|+\ell$ can look like. All proofs use only pair-disjointness, a budget, and counting; none uses optimality of the solution. The lemmas \emph{do not} give an FPT algorithm in $\ell+w$ and no hardness; see Section~\ref{sec:why} for what they do and do not explain, and Section~\ref{sec:open}.

\paragraph{Conventions.} Here a pair is \emph{hard} if it is not packed (in the language of the previous sections: \emph{unpacked}); a \emph{hard edge} is a hard pair in $E$ and a \emph{hard non-edge} a hard pair outside $E$. For a packing $H$ let $c=c(H)$ be the maximum number of members of $H$ through a vertex, $D=D(H)=\max_v|Q(v)|$ the maximum number of packed partners, and $w=\tw(I(H))$ (with $w=0$ if $I(H)$ has no edge). A triple $(G,H,\ell)$ is a \emph{YES instance} if $\OPT(G)\le|H|+\ell$. A \emph{family} $F$ of induced $P_3$'s is \emph{hard over $H$} if the sets $\pairs(P)$, $P\in F$, are pairwise disjoint and disjoint from the packed pairs of $H$; then $H\cup F$ is a packing with $|H|+|F|$ members.

\begin{lemma}[Degrees]\label{lem:deg}
For every vertex $v$, $|Q(v)|=2\,|\{P\in H:v\in P\}|$. Hence $D=2c\le 2(w+1)$.
\end{lemma}
\begin{proof}
Each member $P\ni v$ contributes the two packed pairs of $P$ at $v$, i.e.\ two partners. Two members through $v$ cannot have a common partner $x\ne v$, since they would then share the two vertices $v,x$, hence a pair. So the partner sets are disjoint and $|Q(v)|=2|\{P\ni v\}|$. The inequality $c\le w+1$ is stated in Section~\ref{sec:prelim}.
\end{proof}

\subsection{Maximalisation and the hard-edge lemma}
\begin{lemma}[Hard families]\label{lem:fam}
Let $H$ be any packing, $(G,H,\ell)$ a YES instance and $F$ a family hard over $H$. Then $|F|\le\ell$. In general $\OPT\ge|H|+|F|$.
\end{lemma}
\begin{proof}
$H\cup F$ is a modification-disjoint packing of $|H|+|F|$ induced $P_3$'s. Lemma~\ref{lem:decomp}, applied to $H\cup F$ and an optimum solution, gives $\OPT\ge|H|+|F|$. In a YES instance $\OPT\le|H|+\ell$.
\end{proof}

\begin{lemma}[MAX: maximalisation]\label{lem:MAX}
Let $H$ be any packing and $M$ a maximal family hard over $H$ (inclusion-maximal; it is found greedily in polynomial time). Put $H'=H\cup M$, $\ell'=\ell-|M|$. Then:
\begin{enumerate}
\item[(a)] $\OPT\ge|H'|$. In particular $|M|>\ell$ implies that $(G,H,\ell)$ is a NO instance.
\item[(b)] $\OPT\le|H|+\ell$ iff ($|M|\le\ell$ and $\OPT\le|H'|+\ell'$). So $(G,H,\ell)$ and $(G,H',\ell')$ are equivalent whenever $\ell'\ge0$.
\item[(c)] $H'$ is a packing and is pair-maximal: no induced $P_3$ of $G$ has all three pairs unpacked with respect to $H'$. (``Hard'' now means hard with respect to $H'$; the pairs of $M$ have become packed.)
\item[(d)] $\tw(I(H'))\le\tw(I(H))+|M|$, hence $w'\le w+|M|\le w+\ell$.
\item[(e)] $c'\le\min\{c+|M|,\,w'+1\}$ and $D'=2c'\le D+2|M|$.
\end{enumerate}
A tree decomposition of $I(H')$ of width at most $w+|M|$ is obtained in polynomial time from one of $I(H)$.
\end{lemma}
\begin{proof}
(a) is Lemma~\ref{lem:fam} (first sentence: $\OPT\ge|H|+|M|$). (b) If $\OPT\le|H|+\ell$ then $|M|\le\ell$ by (a), and $|H'|+\ell'=|H|+\ell$. Conversely $|M|\le\ell$ and $\OPT\le|H'|+\ell'=|H|+\ell$. So (b) is arithmetic once (a) is known.

(c) The members of $M$ have pairwise disjoint pair sets, disjoint from those of $H$, so $H'$ is a packing. Suppose an induced $P_3$ $Q$ has all its pairs hard with respect to $H'$. Then $\pairs(Q)$ is disjoint from $\pairs(H)$ and from $\pairs(P)$ for all $P\in M$, so $M\cup\{Q\}$ is a family hard over $H$ ($Q\notin M$ because the pairs of $Q$ are not in $\pairs(M)$), contradicting maximality of $M$.

(d) Let $(T,\{B_t\})$ be a tree decomposition of $I(H)$ of width $w$ (if $H=\emptyset$ take a single empty bag). The graph $I(H')$ has node set $H\cup M$ and $I(H')[H]=I(H)$. Put $B'_t=B_t\cup M$. Every node of $H$ still occupies a connected subtree. Every node of $M$ occupies all of $T$, which is connected. An edge of $I(H')$ with both ends in $H$ is covered as before; an edge with an end $P\in M$ is covered by any bag containing its other end (every node of $H$ lies in some bag; if both ends are in $M$ every bag covers it). Bags have at most $w+1+|M|$ nodes, so the width is at most $w+|M|$.

(e) At most $|M|$ members of $M$ pass through a vertex, so $c'\le c+|M|$; $c'\le w'+1$ is Section~\ref{sec:prelim} for $H'$. By Lemma~\ref{lem:deg}, $D'=2c'\le 2c+2|M|=D+2|M|$.
\end{proof}
The bound in (d) cannot be improved in general.
\begin{example}[Tightness of (d)]\label{ex:maxtight}
Let $G$ be the star with centre $v$ and leaves $p,q,a_1,b_1,\dots,a_m,b_m$, $H=\{(p,v,q)\}$ and $M=\{(a_i,v,b_i):1\le i\le m\}$. The members of $M$ are induced $P_3$'s with pairwise disjoint pair sets, disjoint from $\pairs(H)$. $M$ is maximal: every induced $P_3$ of a star is $(x,v,y)$ with leaves $x,y$, and every pair $xv$ is already packed. $I(H)$ is one node, $\tw=0$; $I(H\cup M)=K_{m+1}$, $\tw=m=0+|M|$.
\end{example}
After applying Lemma~\ref{lem:MAX}, all the other lemmas of this section have to be applied to the primed instance $(G,H',\ell')$ with parameters $w'\le w+|M|$, $c'\le c+|M|$, $D'\le D+2|M|$ and $\ell'=\ell-|M|$. This is the same normalisation as Remark~\ref{rem:max}, with the cost in treewidth made explicit.

\begin{lemma}[HE: hard edges close up]\label{lem:HE}
Let $H$ be pair-maximal and let $ab,bc$ be hard edges ($a,b,c$ distinct). Then $ac\in E$ or $ac$ is packed.
\end{lemma}
\begin{proof}
If $ac\notin E$ and $ac$ were hard, then $a\,b\,c$ would be an induced $P_3$ with all three pairs unpacked, contradicting pair-maximality. So $ac\in E$ or $ac$ is packed.
\end{proof}
No YES hypothesis is needed. Without pair-maximality the lemma fails ($H=\emptyset$ and any induced $P_3$). For an arbitrary packing $H$ one applies it to $H'=H\cup M$ of Lemma~\ref{lem:MAX}(c), with ``hard'' meaning hard with respect to $H'$.

\subsection{Hubs}
\begin{lemma}[HUB$^+$]\label{lem:HUB}
Let $H$ be any packing, $(G,H,\ell)$ a YES instance, $u$ a vertex and $N_h(u)=\{x:ux\in E,\ ux\text{ not packed}\}$. Let $\Gamma_u$ be the graph on $N_h(u)$ whose edges are the pairs $\{x,y\}$ with $xy\notin E$ and $xy$ not packed (hard non-edges). Then:
\begin{enumerate}
\item[(a)] $\nu(\Gamma_u)\le\ell$ (matching number);
\item[(b)] $\Gamma_u$ has a vertex cover $T$ with $|T|\le2\ell$, and every pair of $N_h(u)\setminus T$ is an edge of $G$ or is packed;
\item[(c)] $\alpha(G[N_h(u)])\le 2\ell+c+1$ where $c=c(H)$.
\end{enumerate}
\end{lemma}
\begin{proof}
(a) Let $x_1y_1,\dots,x_ty_t$ be a matching of $\Gamma_u$ and $P_i=(x_i,u,y_i)$. Since $ux_i,uy_i\in E$ and $x_iy_i\notin E$, $P_i$ is an induced $P_3$. Its three pairs $ux_i,uy_i,x_iy_i$ are hard. The $2t$ vertices $x_i,y_i$ are distinct and different from $u$, so the pairs $ux_i,uy_i$ ($i\le t$) are pairwise distinct, and the pairs $x_iy_i$ are pairwise distinct and do not contain $u$. So $\{P_1,\dots,P_t\}$ is hard over $H$ and $t\le\ell$ by Lemma~\ref{lem:fam}.

(b) The vertex set of a maximum matching of $\Gamma_u$ is a vertex cover (by maximality) of size $2\nu\le2\ell$. A pair of $N_h(u)\setminus T$ is not an edge of $\Gamma_u$, i.e.\ it is an edge of $G$ or packed.

(c) Let $J$ be independent in $G[N_h(u)]$. By (b), every pair of $J\setminus T$ is a non-edge of $G$ and not an edge of $\Gamma_u$, hence packed. Put $J'=J\setminus T$. A packed pair of $J'$ lies in a member $P$ of $H$ and is a non-edge, so it is the pair of the two ends of $P$ (the other two pairs of $P$ are edges); in particular $|V(P)\cap J'|\le2$. If $J'\ne\emptyset$ fix $x\in J'$; every $y\in J'\setminus\{x\}$ gives a member $P_y\ni x,y$, and $P_y\ne P_{y'}$ for $y\ne y'$ because $|V(P_y)\cap J'|\le2$. So $|J'|-1\le c$. Hence $|J|\le|T|+|J'|\le2\ell+c+1$.
\end{proof}

\begin{corollary}[HUB, hard-independent sets]\label{cor:HUB}
In the setting of Lemma~\ref{lem:HUB}, if $X\subseteq N_h(u)$ and every pair of $X$ is a hard non-edge, then $|X|\le2\ell+1$.
\end{corollary}
\begin{proof}
$X$ is a clique of $\Gamma_u$, so it contains a matching of size $\lfloor|X|/2\rfloor$, and $\lfloor|X|/2\rfloor\le\ell$ by Lemma~\ref{lem:HUB}(a).
\end{proof}
\begin{remark}[HUB$^+$ versus an independence-number bound]\label{rem:hubalpha}
(1) Bounding $\alpha(G[N_h(u)])$ through the corollary (greedy choice of a subset of an independent set that is also free of packed pairs) gives only $(c+1)(2\ell+1)=(D/2+1)(2\ell+1)$, which is weaker than (c). (2) The matching form (a) is the relevant statement for Vertex Cover on the non-adjacency graph of $N(u)$: let $N_h(u)$ be two disjoint cliques of $G$ of size $m$ with no packed pair between them. Then $\alpha=2$, but $(a_i,u,b_i)$, $a_i,b_i$ from different cliques, are $m$ pair-disjoint induced $P_3$'s, and $m\le\ell$ in a YES instance. Large independent sets are not the obstruction; large matchings of $\Gamma_u$ are. (3) For a general $H$ the lemma is applied to $H'$ of Lemma~\ref{lem:MAX} with $c',\ell'$.
\end{remark}

\subsection{Pieces and small merges}
Fix a packing $H$ and a solution $S$ ($G\sd S$ is a cluster graph) with $|S|\le|H|+\ell$. Let $U=S\setminus\bigcup_i\pairs(P_i)$ be the unpacked part and $u=|U|$; by Lemma~\ref{lem:decomp}, $u\le|S|-|H|\le\ell$. Put $F=U\cap E$ (deleted hard edges) and $I=U\setminus E$ (inserted hard pairs). The \emph{pieces of $S$} are the connected components of the graph $(V,K)$, $K=\{\text{hard edges of }G\}\setminus F$ (kept hard edges); isolated vertices are singleton pieces. \emph{Pieces depend on $S$, not only on $(G,H)$.} Every edge of $K$ is an edge of $G\sd S$, so every piece lies inside one cluster of $G\sd S$.

\begin{lemma}[SM: small merge]\label{lem:SM}
Let $S$, $u$, $I$, pieces as above, $D\ge\max_v|Q(v)|$ (e.g.\ $D=D(H)$), and let $C$ be a cluster of $G\sd S$.
\begin{enumerate}
\item[(i)] If $A\ne B$ are pieces contained in $C$, then every pair $ab$ ($a\in A,b\in B$) is packed or belongs to $I$.
\item[(ii)] For such $A,B$ with $|A|\ge|B|$: $|B|\,(|A|-D)\le|I\cap(A\times B)|\le u$, hence $|A|\le D+u\le D+\ell$.
\item[(iii)] If $C$ is the union of at least two pieces, then $|C|\le\max\{D+u+1,\;2D,\;2u\}$.
\end{enumerate}
\end{lemma}
\begin{proof}
(i) Since $a,b\in C$, $ab$ is an edge of $G\sd S$, i.e.\ $ab\in E$ xor $ab\in S$. Suppose $ab$ is hard. If $ab\in E$ then $ab\notin S$; as $F\subseteq S$, $ab\notin F$, so $ab\in K$ and $a,b$ lie in one piece, a contradiction. So $ab\notin E$, hence $ab\in S$, and as $ab$ is hard, $ab\in U\setminus E=I$.

(ii) Each vertex of $B$ has at most $D$ packed partners, so at most $D|B|$ pairs of $A\times B$ are packed. By (i) all other pairs of $A\times B$, at least $|A||B|-D|B|$, lie in $I$, and $|I|\le|U|=u$. If $|A|\ge D+u+1$ then $|B|(|A|-D)\ge u+1$ because $|B|\ge1$, a contradiction; so $|A|\le D+u$.

(iii) Let $A$ be a largest piece in $C$, $a=|A|$ and $r=|C|-a\ge1$. By (i), every pair $\{x,y\}$ with $x,y\in C$ in different pieces is packed or in $I$.

\emph{Case $a>D$.} Each $x\in C\setminus A$ has at most $D$ packed partners, so at least $a-D\ge1$ vertices $y\in A$ give a pair $xy$ in $I$. These pairs are distinct for distinct $(x,y)$, so $r(a-D)\le|I|\le u$. Then $r\le u$ and $a\le D+u/r$, so $|C|=a+r\le D+u/r+r$. The function $r\mapsto r+u/r$ is convex on $[1,u]$ and takes the value $u+1$ at both ends, so $|C|\le D+u+1$. (The case does not occur if $u=0$.)

\emph{Case $a\le D$.} All pieces of $C$ have at most $D$ vertices. Each $x\in C$ has at least $|C|-D$ vertices of $C$ outside its piece, at most $D$ of which are packed partners; so $x$ lies in at least $|C|-2D$ pairs $\{x,y\}$ of $I$ with $y\in C$. Summing over $x$ counts each pair twice: $2|I|\ge|C|(|C|-2D)$. If $|C|>2D$, then $|C|-2D\ge1$ and $|C|\le2|I|\le2u$. So $|C|\le\max\{2D,2u\}$.
\end{proof}
\begin{remark}[Qualifiers and relation to the original bound]\label{rem:SMqual}
(1) The qualifier ``pieces sharing a cluster'' is essential: a single piece may be an entire cluster of unbounded size, and a statement ``pieces have size at most $D+\ell$'' is false. (2) The bound (iii) implies the weaker bound $|C|\le\max\{2D+\ell,2\ell\}$ with which the lemma first appeared in our notes: $2u\le2\ell$ and $2D\le2D+\ell$ are clear; $D+u+1\le2D+\ell$ if $D\ge1$; if $D=0$ and $u\ge1$ then $u+1\le2u\le2\ell$; and if $D=u=0$ no cluster with two pieces exists (the two cases of the proof are both impossible). (3) The bound (iii) is attained in some small instances of our tests (slack $0$), for instance two singleton pieces and one inserted pair; we do not claim it is tight in general. (4) Lemma~\ref{lem:SM} says nothing about the \emph{degrees} in a graph of pieces; it bounds the sizes of pieces that share a cluster and the number of cross pairs that are not packed.
\end{remark}

\subsection{The soft graph has bounded treewidth}
Let $W=(V,\{\text{packed pairs}\})$ be the \emph{soft graph}.
\begin{proposition}[Treewidth of the soft graph]\label{prop:soft}
$\tw(W)\le 3(w+1)-1=3w+2$ if $H\ne\emptyset$; $W$ is edgeless if $H=\emptyset$. The bound is attained for $w=0$ ($H$ a single $P_3$, $W$ a triangle).
\end{proposition}
\begin{proof}
Let $(T,\{B_t\})$ be a tree decomposition of $I(H)$ of width $w$ and $T_P=\{t:P\in B_t\}$, a subtree of $T$ for each $P\in H$. Put $B'_t=\bigcup_{P\in B_t}V(P)$, so $|B'_t|\le3(w+1)$. For every vertex $x\notin V(H)$ add a new leaf bag $\{x\}$ to an arbitrary node. \emph{Edges:} a packed pair lies in $V(P)$ for some $P\in H$, and $P\in B_t$ for some $t$, so the pair lies in $B'_t$. \emph{Connectivity:} for $x\in V(H)$ the nodes $t$ with $x\in B'_t$ form $\bigcup_{P\ni x}T_P$. The members through $x$ pairwise share $x$, so they are pairwise adjacent in $I(H)$ and, the edge being covered by a bag, their subtrees $T_P$ pairwise intersect. A union of connected subtrees that pairwise intersect is connected (two points of the union lie in subtrees $T_P,T_{P'}$ that meet; join them through a common node). The new leaf bags are single bags of size $1$. So this is a tree decomposition of $W$ of width at most $3(w+1)-1$.
\end{proof}
Computationally, for $w\ge1$ the largest value of $\tw(W)$ that we saw on small instances is much smaller than $3w+2$; we do not claim the bound is tight beyond $w=0$.

\subsection{Computer verification}
All statements of this section were checked by brute force with code written independently of the other scripts (\texttt{code/w5\_packing\_lemmas.py}, all set partitions, $n\le8$, 4000 random graphs: random, noisy clusters, hubs; for each a random partial packing $H$ and its maximalisation $H\cup M$): Lemma~\ref{lem:deg} ($D=2c$), Lemma~\ref{lem:MAX}(a)--(e), Example~\ref{ex:maxtight} for $m\le5$, Lemma~\ref{lem:HE}, Lemma~\ref{lem:HUB} and Corollary~\ref{cor:HUB} for every vertex and every $\ell$ with $\OPT\le|H|+\ell\le|H|+3$, Lemma~\ref{lem:SM} on every solution of cost at most $|H|+3$ (about $5.7\cdot10^5$ clusters with two or more pieces, including slack $0$ in (iii) on some cases), and Proposition~\ref{prop:soft} (exact treewidth). No violation; the variants $\alpha\le2\ell+c$ and $|C|\le\max\{D+u,2D,2u\}$ are falsified by the same script, so the checks are not vacuous. This is evidence for the statements; the proofs above are the justification.

\subsection{Why simple hardness encodings fail}\label{sec:why}
\emph{This subsection is a heuristic discussion, not a theorem.} It explains why three natural ways of building W[1]-hardness for $\ell+w$ do not work as stated; it does not exclude other gadgets and it does not prove fixed-parameter tractability, which remains open (Open Problem~\ref{op:main}). A reduction from a problem with parameter $k$ must keep $\ell+w\le g(k)$; grids of side $k$ inside $I(H)$ are therefore allowed, but structures that grow with the instance size are not.

\paragraph{Hub of deletions.} A hub $u$ is adjacent to vertices $x_1,\dots,x_m$ and the instance asks to keep a subset of $N(u)$ forming a clique and delete the rest, i.e.\ Vertex Cover on the non-adjacency graph of $N(u)$. Lemma~\ref{lem:HUB} shows what survives in a YES instance: the graph of \emph{hard} non-adjacent pairs inside $N_h(u)$ has a matching of size at most $\ell$ and a vertex cover of size at most $2\ell$, so at most $2\ell$ hub neighbours carry hard conflicts; the other non-adjacencies are adjacencies or packed pairs, and a vertex has at most $2c$ of the latter. The relevant obstruction is a large matching (Remark~\ref{rem:hubalpha}(2)), not a large independent set, and any such matching immediately forces a NO answer. Non-adjacencies realised by packed pairs are local, with degree at most $2c\le2(w+1)$ in the soft graph. This is why Vertex Cover through a hub cannot encode an unbounded source graph; it does not address hubs whose conflicts are realised by packed pairs; a vertex has at most $2c$ of these, so many of them force large $c$, hence large $w$.

\paragraph{Hub of insertions.} A hub $u$ with classes $V_1,\dots,V_k$ chooses, by inserting $uv_{i,j_i}$, one vertex per class; the selected vertices land in one cluster $C$ and compatibility is encoded by adjacency of the chosen vertices. By Lemma~\ref{lem:SM}(i), two selected vertices in different pieces of $C$ form a packed pair or an inserted pair (at most $\ell$ of these); adjacency by a hard edge that is not deleted puts them in the same piece, and Lemma~\ref{lem:SM}(ii) bounds the size of merged pieces by $D+\ell$. So the compatibility relation between selected vertices in different pieces can only be realised by packed pairs, which form the soft graph, of treewidth at most $3w+2$ (Proposition~\ref{prop:soft}), or by a bounded number of inserted pairs. The caveats: Lemma~\ref{lem:SM} bounds sizes of pieces, not the degree of a compatibility graph; at most $\ell$ hard edges may be deleted, which a gadget could exploit; and pieces depend on the chosen solution. The argument shows that the hard graph cannot serve as a channel for $[n]$-valued information between selected vertices, nothing more.

\paragraph{Grid tiling.} A tile vertex $t_{ab}$ is joined by hard edges to a row vertex $x_a$ and a column vertex $y_b$ ($x_ay_b$ a non-edge). The triples $(x_a,t_{ab},y_b)$ are induced $P_3$'s; if their pairs are hard they are pairwise pair-disjoint (their pairs are $x_at_{ab}$, $t_{ab}y_b$, $x_ay_b$, all different for different $(a,b)$), so by Lemma~\ref{lem:fam} at most $\ell$ of them can be hard, whereas a grid-tiling encoding needs about $n^2$ tiles and $n$ is not bounded by $g(k)$. So all but at most $\ell$ of them are packed; then some row vertex $x_a$ lies in at least $n-\ell$ members (if every row has $n$ tiles), $c\ge n-\ell$, hence $w\ge n-\ell-1$, again unbounded in $k$. Mixed and replicated gadgets (several copies of row and column vertices, tiles realised by chains) were not excluded by a proof; the examination in our notes found no variant that carries information on $[n]$ across a tree decomposition of bounded width, which is evidence only.

\paragraph{What remains open.} The lemmas above make every ``simple'' source of unbounded information in $I(H)$-width-bounded instances expensive. They do not bound the candidate positions of inserted pairs between pieces. Sections~\ref{sec:locality} and~\ref{sec:dproute} go further: the inserted pairs of an optimal solution are local for type T1, in particular for merges of pieces with no soft pair between them, but not in general (Proposition~\ref{prop:cex}), and given one endpoint of an edited hard pair the other has a bounded candidate set (Corollary~\ref{cor:cand}). The FPT question for $\ell+w$ stays an open problem.
